\section{Task 1: Text Generation Models}
\label{sec:task1}

Six models are trained and evaluated in Task 1. They form a 3$\times$2 grid:
three architectures (Markov Chain, RNN, LSTM) applied at two tokenization levels
(character, word).

\subsection{Markov Chain}

The Markov Chain implementation (\texttt{models/markov\_chain.py}) builds a
transition table from consecutive n-gram sequences. No neural network is involved.

\begin{itemize}
    \item \textbf{Character-level} uses order $n = 5$: the next character is
          sampled from the empirical distribution of characters that followed the
          same 5-character context in the training corpus.
    \item \textbf{Word-level} uses order $n = 2$: the next word is sampled given
          the previous two words.
\end{itemize}

Generation simply looks up the current state in the transition table and samples
proportionally to observed frequencies. If a state has never been seen, the model
restarts from a random known starting n-gram.

\subsection{Recurrent Neural Network (RNN)}

The RNN model (\texttt{models/rnn\_model.py}) processes tokens sequentially via
a 2-layer vanilla RNN:

\begin{enumerate}
    \item An embedding layer maps each token index to a dense vector.
    \item Two RNN layers with \texttt{tanh} activation update the hidden state.
    \item A linear output layer projects the hidden state to a vocabulary-size
          logit vector.
\end{enumerate}

Training uses cross-entropy loss with gradient clipping (max norm $= 1.0$) and
the Adam optimiser. Temperature sampling ($T = 0.8$) is used during generation
to balance diversity and coherence.

\subsection{Long Short-Term Memory (LSTM)}

The LSTM model (\texttt{models/lstm\_model.py}) has the same architecture as the
RNN but replaces the RNN cells with LSTM cells, adding a cell state and three
gates. The gating mechanism allows gradients to flow without vanishing, enabling
the model to capture patterns across many more timesteps.

\subsection{Training Configuration}

All neural models (RNN and LSTM) are trained for \textbf{20 epochs} with:
\begin{itemize}
    \item Optimiser: Adam
    \item Gradient clipping: max norm $= 1.0$
    \item Batch size: 64
    \item Sequence length: 100 tokens (char), 20 tokens (word)
    \item Generation temperature: $T = 0.8$
\end{itemize}
